\section{Introduction}
\label{sec:introduction}
Quantum computers promise to solve problems of practical importance that are intractable for classical computers, such as simulating quantum systems~\cite{lloyd1996universal} and integer factoring~\cite{Shor1994Factoring}.
Existing quantum computers are highly susceptible to noise, which prevents them from executing complex algorithms such as Shor's factoring algorithm.
To execute such complex quantum algorithms, we need fault-tolerant quantum computers (FTQC)~\cite{shor1996fault}, which rely on quantum error correction (QEC)~\cite{shor1995scheme, Steane1996ErrorCorrecting}.

In FTQC, each \textit{logical gate} must be implemented in a fault-tolerant manner.
While fault-tolerant implementations of Clifford gates are relatively inexpensive, Clifford gates alone are not sufficient for universal quantum computation~\cite{gottesman1998heisenberg}.
Implementing a non-Clifford gate such as the $T$ gate requires a special quantum state, a \textit{magic state}, which is obtained through a process known as \textit{magic state distillation}~\cite{Bravyi2005}.
Magic state distillation is resource-intensive and remains a major bottleneck in fault-tolerant quantum computation.
Consequently, many distillation protocols have been proposed~\cite{Bravyi2005, knill1998resilient, bravyi2012magic, meier2012magic, Goto2016Minimizing, Campbell2017Unified, Gidney2019EfficientMagicState, Litinski2019GameOfSurfaceCodes, Litinski2019magicstate, chamberland2020very, itogawa2025efficient, Hirano2024Leveraging, Gidney2024MagicStateCultivation, vaknin2025magic, chen2025efficient, sahay2025fold, claes2025cultivating}.

Bravyi~\textit{et al.}~\cite{Bravyi2005} proposed the 15-to-1 distillation protocol, which exploits $T$'s transversality on the $[[15, 1, 3]]$ code.
Earlier, predating the term ``magic state distillation'', Knill~\textit{et al.}~\cite{knill1998resilient} proposed a protocol based on the Hadamard test of the $H$ gate, relying on $H$'s transversality on the $[[7, 1, 3]]$ code.
These studies formulated magic state distillation as a quantum algorithm built from ideal Clifford gates with noisy inputs, and this formulation was followed by subsequent studies~\cite{bravyi2012magic, meier2012magic, Campbell2017Unified, Gidney2019EfficientMagicState, Litinski2019GameOfSurfaceCodes}.
Running such algorithms on FTQC requires protecting each gate against noise using QEC, incurring a significant cost.

Goto~\cite{Goto2016Minimizing} proposed a technique that we call \textit{physical-level distillation}.
In contrast to the above formulation, the distillation circuit operates with noisy Clifford gates.
Specifically, the circuit detects any single gate error, thereby achieving a logical error probability of $\mathcal{O}(p_\mathrm{phys}^2)$, where $p_\mathrm{phys}$ is the physical error probability.
To this end, the distillation protocol consists of a Knill-style Hadamard test with a flag qubit.
Its physical-gate implementation significantly reduces the distillation overhead.

Goto's protocol requires all-to-all connectivity.
Chamberland~\textit{et al.}~\cite{chamberland2020very} and Itogawa~\textit{et al.}~\cite{itogawa2025efficient} proposed physical-level distillation protocols that work with planar connectivity.
The latter, known as zero-level distillation, operates on the simpler square-grid connectivity and yields a magic state encoded in a surface code~\cite{Kitaev2003, bravyi2007measurement}, achieving a logical error probability of $10^{-4}$ for a realistic $p_\mathrm{phys} = 10^{-3}$ setting.
Gidney~\textit{et al.}~\cite{Gidney2024MagicStateCultivation} refined these works and proposed \textit{magic state cultivation} (MSC), which works with the square-grid connectivity.
This protocol extensively utilizes postselection to suppress errors, thereby producing high-fidelity magic states.
As a result, it yields a magic state with a logical error probability of $2 \times 10^{-9}$ at $p_\mathrm{phys} = 10^{-3}$.
MSC prepares a magic state encoded in a 2D color code~\cite{bombin2006topological}, which is then transformed into the grafted code, another QEC code compatible with matching QEC decoders.
However, due to its high-weight stabilizers, the grafted code is less space-efficient than the rotated surface code~\cite{bombin2007optimal} with a comparable logical error probability, incurring substantial spacetime overhead.

Although these physical-level distillation protocols significantly reduce the distillation overhead, further advances are desirable, particularly when many magic states are consumed simultaneously during the parallel execution of multiple $T$ gates~\cite{Beverland2022Assessing, hirano2025locality}.
While this has motivated the development of various MSC-based protocols~\cite{vaknin2025magic, chen2025efficient, sahay2025fold, claes2025cultivating}, they relax the connectivity restriction that MSC respects.
This restriction is particularly important for superconducting devices.
Additionally, although platforms such as neutral atoms allow flexible connectivity, qubit shuttling incurs temporal and idling costs; thus, maintaining local operations on a fixed lattice can offer temporal advantages.
Therefore, there is a need for a more efficient magic state distillation protocol compatible with the square-grid connectivity.

In this study, we propose a novel MSC-based protocol that significantly reduces the spacetime overhead of distillation without relaxing the hardware connectivity requirements.
We avoid using the grafted code by employing lattice surgery~\cite{Horsman2012, poulsen2017fault}, as in the original zero-level distillation~\cite{itogawa2025efficient}, to transfer the magic state from the color code to the rotated surface code.
Lattice surgery and the rotated surface code allow us to employ two techniques to further reduce the distillation overhead:
(1) utilizing surface code expansion~\cite{li2015magic} to minimize the qubit count in the early stages of distillation, and (2) predicting the logical error probability using a simple lookup table, allowing early rejection to reduce the average spacetime overhead.

In addition to efficiency, our protocol provides simplicity.
The grafted code is complex, and avoiding its implementation is appealing to practitioners.
The rotated surface code and lattice surgery are simpler, already supported by existing compilers~\cite{Watkins2023High} and instruction sets, and impose little additional burden on control hardware.
Thus, our protocol can leverage existing software and hardware infrastructures, making it more practical for near-term implementations.

We conducted numerical simulations using Stim~\cite{gidney2021stim} to evaluate the performance of our protocol.
With a color code distance of 3, our protocol achieves a logical error probability for the resulting magic state comparable to that of MSC, while requiring about half the spacetime overhead.
Specifically, at $p_\mathrm{phys} = 10^{-3}$, our protocol achieves a logical error probability of $1.6 \times 10^{-6}$ for $S\ket{+}$ distillation.
This is comparable to the $1.2 \times 10^{-6}$ achieved by MSC with color code distance 3, and corresponds to a logical error probability of $3 \times 10^{-6}$ for $T\ket{+}$ distillation.
For logical error probabilities in the range $[10^{-6}, 10^{-5}]$, our protocol requires less than half the spacetime overhead needed for MSC.
These results demonstrate that our approach offers a practical path toward more efficient magic state distillation in fault-tolerant quantum computation, especially for megaquop use cases~\cite{preskill2025beyond} (computation involving $\sim 10^{6}$ logical gate operations) and early fault-tolerant devices~\cite{akahoshi2024partially,toshio2025practical, akahoshi2024compilation} (systems with $10^3$--$10^5$ physical qubits).

Our contributions are summarized as follows:
\begin{enumerate}
 \item We demonstrate that, although the original MSC work found lattice-surgery-based constructions unsuitable, lattice surgery can in fact be used within MSC-style distillation to realize a simpler and more efficient distillation protocol.
 \item We define a concrete protocol, magic state cultivation with lattice surgery (MSC-LS), an efficient and simple distillation protocol using lattice surgery and surface code expansion. This protocol is compatible with the square-grid connectivity and avoids the use of the grafted code.
 \item We introduce an early-rejection mechanism based on a lookup table of intermediate error syndrome patterns. This mechanism reduces the expected spacetime overhead by aborting attempts that are unlikely to pass the final postselection.
 \item We quantitatively evaluate MSC-LS and show that it achieves logical error probabilities comparable to those of MSC while reducing the expected spacetime overhead by roughly 50\%.
\end{enumerate}

The rest of this paper is structured as follows.
In \autoref{sec:preliminary}, we provide basic definitions and prior work essential for understanding our proposal.
Specifically, we define the detector error model and provide a summary of the complementary gap and magic state cultivation.
In \autoref{sec:msc-bottleneck}, we discuss bottlenecks in magic state cultivation that motivate our proposal.
In \autoref{sec:msc-ls}, we detail our proposal: magic state cultivation with lattice surgery.
In \autoref{sec:performance-evaluation}, we present a performance evaluation of the protocol, including the numerical simulation settings and results.
Finally, \autoref{sec:conclusion} summarizes our findings.